\documentclass[11pt]{article}
\usepackage{amsmath}
\usepackage{geometry}
\usepackage{booktabs}
\usepackage{xcolor}
\usepackage{listings}

\geometry{margin=1in}

\lstset{
    basicstyle=\ttfamily\small,
    breaklines=true,
    columns=fullflexible,
    keepspaces=true,
    showstringspaces=false,
    frame=single,
    backgroundcolor=\color{gray!10}
}

\title{A Brief Primer on Zsh Commands}
\author{For macOS Terminal}
\date{}

\begin{document}
\maketitle

\tableofcontents
\newpage

% ---------------------------------------------------------------
\section{Introduction}

Zsh (Z shell) is the default shell on macOS since Catalina (10.15). It's backwards-compatible with bash while offering enhanced features like better tab completion, glob patterns, and themeable prompts.

To check your current shell:
\begin{lstlisting}
echo $SHELL
\end{lstlisting}

% ---------------------------------------------------------------
\section{Navigation and File Management}

\subsection{Basic Navigation}

\begin{center}
\begin{tabular}{ll}
\toprule
Command & Description \\
\midrule
\texttt{pwd} & Print working directory \\
\texttt{ls} & List files and directories \\
\texttt{ls -la} & List all files (including hidden) with details \\
\texttt{cd \textasciitilde} & Change to home directory \\
\texttt{cd ..} & Go up one directory \\
\texttt{cd -} & Go to previous directory \\
\bottomrule
\end{tabular}
\end{center}

\subsection{File Operations}

\begin{lstlisting}
# Create files and directories
touch filename.txt          # Create empty file
mkdir dirname               # Create directory
mkdir -p path/to/nested/dir # Create nested directories

# Copy, move, delete
cp file.txt copy.txt        # Copy file
cp -r dir/ newdir/          # Copy directory recursively
mv oldname.txt newname.txt  # Rename/move file
rm file.txt                 # Remove file
rm -r dirname/              # Remove directory recursively
rm -rf dirname/             # Force remove (use carefully!)
\end{lstlisting}

% ---------------------------------------------------------------
\section{Viewing and Editing Files}

\begin{center}
\begin{tabular}{ll}
\toprule
Command & Description \\
\midrule
\texttt{cat file.txt} & Display entire file \\
\texttt{less file.txt} & View file page by page (q to quit) \\
\texttt{head file.txt} & Show first 10 lines \\
\texttt{tail file.txt} & Show last 10 lines \\
\texttt{tail -f log.txt} & Follow file updates in real-time \\
\texttt{nano file.txt} & Edit file with nano editor \\
\texttt{vim file.txt} & Edit file with vim editor \\
\bottomrule
\end{tabular}
\end{center}

% ---------------------------------------------------------------
\section{Search and Find}

\begin{lstlisting}
# Find files
find . -name "*.txt"           # Find all .txt files
find . -type d -name "dirname" # Find directories

# Search within files
grep "search term" file.txt    # Search for text in file
grep -r "term" .               # Recursive search in directory
grep -i "term" file.txt        # Case-insensitive search
grep -n "term" file.txt        # Show line numbers
\end{lstlisting}

% ---------------------------------------------------------------
\section{Zsh-Specific Features}

\subsection{Advanced Globbing}

Zsh offers powerful pattern matching:

\begin{lstlisting}
# Extended globbing
ls **/*.txt                # Recursively find all .txt files
ls *.{jpg,png}             # Files ending in .jpg or .png
ls file<1-100>.txt         # file1.txt through file100.txt

# Qualifiers
ls *(.)                    # Only files
ls *(/)                    # Only directories
ls *(m-7)                  # Modified in last 7 days
ls *(.L+1m)                # Files larger than 1MB
\end{lstlisting}

\subsection{Directory Stack}

\begin{lstlisting}
# Navigate between directories
pushd /path/to/dir         # Push directory onto stack
popd                       # Pop and return to previous
dirs -v                    # View directory stack
cd ~3                      # Jump to 3rd directory in stack
\end{lstlisting}

\subsection{Tab Completion}

Zsh has intelligent tab completion:
\begin{itemize}
    \item Type partial command + \texttt{Tab} to complete
    \item \texttt{Tab} again to cycle through options
    \item Works for files, directories, commands, and arguments
    \item Case-insensitive by default
\end{itemize}

% ---------------------------------------------------------------
\section{Process Management}

\begin{center}
\begin{tabular}{ll}
\toprule
Command & Description \\
\midrule
\texttt{ps aux} & List all running processes \\
\texttt{top} & Interactive process viewer \\
\texttt{htop} & Enhanced process viewer (if installed) \\
\texttt{kill PID} & Terminate process by ID \\
\texttt{killall name} & Kill all processes by name \\
\texttt{command \&} & Run command in background \\
\texttt{jobs} & List background jobs \\
\texttt{fg} & Bring background job to foreground \\
\texttt{Ctrl+Z} & Suspend current process \\
\texttt{bg} & Resume suspended job in background \\
\bottomrule
\end{tabular}
\end{center}

% ---------------------------------------------------------------
\section{Permissions and Ownership}

\begin{lstlisting}
# View permissions
ls -l file.txt             # Shows: -rw-r--r--

# Change permissions
chmod 755 script.sh        # rwxr-xr-x (executable)
chmod +x script.sh         # Add execute permission
chmod -w file.txt          # Remove write permission

# Change ownership
chown user:group file.txt  # Change owner and group
sudo chown root file.txt   # Change owner to root
\end{lstlisting}

\textbf{Permission notation:}
\begin{itemize}
    \item r (read) = 4
    \item w (write) = 2
    \item x (execute) = 1
    \item 755 = rwxr-xr-x (owner: all, group/others: read+execute)
    \item 644 = rw-r--r-- (owner: read+write, others: read only)
\end{itemize}

% ---------------------------------------------------------------
\section{Input/Output Redirection}

\begin{lstlisting}
# Redirection operators
command > file.txt         # Redirect output to file (overwrite)
command >> file.txt        # Append output to file
command < input.txt        # Read input from file
command 2> error.log       # Redirect errors to file
command &> all.log         # Redirect both output and errors

# Pipes
ls -l | grep ".txt"        # Pipe output to grep
cat file.txt | wc -l       # Count lines
history | tail -20         # Last 20 commands
\end{lstlisting}

% ---------------------------------------------------------------
\section{Variables and Environment}

\begin{lstlisting}
# Variables
NAME="value"               # Set variable (no spaces!)
echo $NAME                 # Use variable
export PATH=$PATH:/new/path # Add to PATH

# Common environment variables
echo $HOME                 # Home directory
echo $USER                 # Current username
echo $SHELL                # Current shell
echo $PATH                 # Executable search path

# Zsh-specific
echo $ZSH_VERSION          # Zsh version
\end{lstlisting}

% ---------------------------------------------------------------
\section{Aliases and Functions}

\subsection{Aliases}

Add to \texttt{\textasciitilde/.zshrc} file:

\begin{lstlisting}
# Common aliases
alias ll='ls -lah'
alias ..='cd ..'
alias ...='cd ../..'
alias gs='git status'
alias gp='git push'
alias gc='git commit'

# After editing .zshrc, reload:
source ~/.zshrc
\end{lstlisting}

\subsection{Functions}

\begin{lstlisting}
# Define function in .zshrc
mkcd() {
    mkdir -p "$1" && cd "$1"
}

# Usage:
mkcd new/nested/directory
\end{lstlisting}

% ---------------------------------------------------------------
\section{History and Shortcuts}

\subsection{Command History}

\begin{lstlisting}
history                    # Show command history
!!                         # Repeat last command
!-2                        # Repeat 2nd-to-last command
!123                       # Repeat command #123
!grep                      # Repeat last grep command
^old^new                   # Replace 'old' with 'new' in last command
\end{lstlisting}

\subsection{Keyboard Shortcuts}

\begin{center}
\begin{tabular}{ll}
\toprule
Shortcut & Action \\
\midrule
\texttt{Ctrl+A} & Go to beginning of line \\
\texttt{Ctrl+E} & Go to end of line \\
\texttt{Ctrl+U} & Clear line before cursor \\
\texttt{Ctrl+K} & Clear line after cursor \\
\texttt{Ctrl+W} & Delete word before cursor \\
\texttt{Ctrl+L} & Clear screen \\
\texttt{Ctrl+R} & Reverse search history \\
\texttt{Ctrl+C} & Cancel current command \\
\texttt{Ctrl+D} & Exit shell / EOF \\
\texttt{Tab} & Auto-complete \\
\bottomrule
\end{tabular}
\end{center}

% ---------------------------------------------------------------
\section{Network and System}

\begin{lstlisting}
# Network
ping google.com            # Test connectivity
curl https://example.com   # Fetch URL content
wget https://file.com/file # Download file (if installed)
ifconfig                   # Network interface info
netstat -an                # Network connections

# System information
uname -a                   # System information
df -h                      # Disk usage
du -sh *                   # Directory sizes
free -h                    # Memory usage (Linux)
uptime                     # System uptime
date                       # Current date/time

# macOS specific
sw_vers                    # macOS version
system_profiler            # Detailed system info
open .                     # Open current directory in Finder
open file.txt              # Open file with default app
pbcopy < file.txt          # Copy file contents to clipboard
pbpaste > file.txt         # Paste clipboard to file
\end{lstlisting}

% ---------------------------------------------------------------
\section{Package Management}

macOS uses Homebrew for package management:

\begin{lstlisting}
# Install Homebrew (if not installed)
/bin/bash -c "$(curl -fsSL https://raw.githubusercontent.com/Homebrew/install/HEAD/install.sh)"

# Common Homebrew commands
brew install package_name  # Install package
brew uninstall package     # Remove package
brew update                # Update Homebrew
brew upgrade               # Upgrade all packages
brew list                  # List installed packages
brew search query          # Search for packages
brew info package          # Package information
\end{lstlisting}

% ---------------------------------------------------------------
\section{Customizing Zsh}

\subsection{Configuration File}

Your Zsh configuration lives in \texttt{\textasciitilde/.zshrc}. Edit with:
\begin{lstlisting}
nano ~/.zshrc
# or
vim ~/.zshrc
\end{lstlisting}

After editing, reload:
\begin{lstlisting}
source ~/.zshrc
\end{lstlisting}

\subsection{Oh My Zsh (Optional)}

Oh My Zsh is a popular framework for managing Zsh configuration:

\begin{lstlisting}
# Install Oh My Zsh
sh -c "$(curl -fsSL https://raw.githubusercontent.com/ohmyzsh/ohmyzsh/master/tools/install.sh)"

# Configuration file becomes: ~/.zshrc
# Themes directory: ~/.oh-my-zsh/themes/
# Plugins directory: ~/.oh-my-zsh/plugins/
\end{lstlisting}

Popular plugins to add to \texttt{.zshrc}:
\begin{lstlisting}
plugins=(git sudo z zsh-autosuggestions zsh-syntax-highlighting)
\end{lstlisting}

% ---------------------------------------------------------------
\section{Quick Reference Table}

\begin{center}
\small
\begin{tabular}{lll}
\toprule
Category & Command & Description \\
\midrule
Navigation & \texttt{cd} & Change directory \\
& \texttt{pwd} & Print working directory \\
& \texttt{ls} & List contents \\
\midrule
Files & \texttt{touch} & Create file \\
& \texttt{cp} & Copy \\
& \texttt{mv} & Move/rename \\
& \texttt{rm} & Remove \\
\midrule
Viewing & \texttt{cat} & Display file \\
& \texttt{less} & Page through file \\
& \texttt{head/tail} & First/last lines \\
\midrule
Search & \texttt{find} & Find files \\
& \texttt{grep} & Search in files \\
\midrule
Help & \texttt{man command} & Manual for command \\
& \texttt{command --help} & Quick help \\
& \texttt{which command} & Find command location \\
\bottomrule
\end{tabular}
\end{center}

% ---------------------------------------------------------------
\section{Tips and Best Practices}

\begin{enumerate}
    \item \textbf{Use tab completion} --- save time and avoid typos
    \item \textbf{Be careful with \texttt{rm -rf}} --- it deletes permanently
    \item \textbf{Use \texttt{man} pages} --- \texttt{man command} shows documentation
    \item \textbf{Test commands with \texttt{echo}} --- verify before execution
    \item \textbf{Use \texttt{Ctrl+R}} --- search command history interactively
    \item \textbf{Create aliases} --- for frequently used commands
    \item \textbf{Comment your scripts} --- use \texttt{\#} for comments
    \item \textbf{Use quotes} --- for filenames with spaces: \texttt{"my file.txt"}
\end{enumerate}

\section{Getting Help}

\begin{lstlisting}
man command                # Full manual page
command --help             # Quick help
command -h                 # Short help (sometimes)
which command              # Find command location
whence -v command          # Zsh: command type and location
type command               # Show how command will be interpreted
\end{lstlisting}

% ---------------------------------------------------------------
\section{Conclusion}

This primer covers the essential Zsh commands for macOS terminal use. For more advanced topics, explore:
\begin{itemize}
    \item Shell scripting
    \item Regular expressions
    \item Process substitution
    \item Zsh parameter expansion
    \item Custom completion functions
\end{itemize}

Practice regularly, and the command line will become second nature!

\end{document}